\documentclass[10pt,a4paper]{amsart}
\usepackage[margin=19mm]{geometry}
\usepackage{amsmath,amssymb,mathtools}
\usepackage[scr=rsfs,cal=euler]{mathalfa}
\usepackage{xfrac}
\usepackage[unicode,hidelinks,bookmarks=false]{hyperref}
\newcommand{\ie}{i.e.\ }
\newcommand{\cf}{cf.\ }
\newcommand{\resp}{respectively\ }
\newcommand{\Subsection}{\subsection}
\providecommand{\nobreakdash}{\nobreak\hbox{-}}
\setlength{\emergencystretch}{2em}
\multlinegap0pt
\usepackage{bm}
\usepackage{bbm}
\usepackage{dsfont}
\usepackage{xspace}
\usepackage{cancel}
\usepackage{mathtools}
\usepackage{amsthm}
\makeatletter
\@ifundefined{thm}{\newtheorem{thm}{Theorem}}{}
\makeatother
\newcommand{\beas}{\begin{eqnarray*}}
\newcommand{\eeas}{\end{eqnarray*}}
\newcommand{\ol}{\overline}
\renewcommand{\epsilon}{\varepsilon}
\title[On the pre-Schwarzian and Schwarzian derivatives of log-harm]{On the pre-Schwarzian and Schwarzian derivatives of log-harmonic mappings}
\author{Raju Biswas, Rajib Mandal}
\date{Source: arXiv:2511.05076v1 (CC BY 4.0)}
\begin{document}
\maketitle

\noindent\emph{Source and licence.} On the pre-Schwarzian and Schwarzian derivatives of log-harmonic mappings. Version arXiv:2511.05076v1 (CC BY 4.0), distributed under the Creative Commons Attribution 4.0 International licence, as recorded in the arXiv licence metadata for this exact version and shown on the abstract page https://arxiv.org/abs/2511.05076v1. Full rights notice: licenses/arxiv-2511.05076-CC-BY-4.0.txt.\par\smallskip

\noindent\textit{Editorial scope.} Counted unit: the Thm whose source label is \texttt{Th11} in the article (source0.tex lines 225--235, source label \texttt{Th11}), delivered together with the article's own complete original proof of it (source0.tex lines 236--276, one \texttt{proof} environment of 41 lines). No mathematics is written, repaired or re-derived: statement and proof are the article's own text line for line. Required context retained uncounted: Th1 (lines 183--192). Every definition, display and label of those lines is retained unchanged. Omitted from this package and not redistributed: 1--182, 193--224, 277--575. Nothing inside a retained excerpt is omitted or altered: every retained body is delivered exactly as the article's lines are written, so no omission receipt of any kind accompanies this package. Citations: the retained excerpts carry no citation command at all, so no bibliography entry of the article is transcribed and no reference is redistributed. Reference adaptation: none was needed and none was made; every reference inside the delivered unit resolves inside it, so no unresolved reference or citation is produced. Printed slips kept literal: no printed word, symbol, hypothesis or number of the counted statement or of its proof was altered. Layout and class: the article's own class is \texttt{amsart} with packages including indentfirst, amssymb, amsmath, amsthm, mathrsfs, cite, graphicx, float; the wrapper is the lane's pinned \texttt{amsart} editorial wrapper at 10pt on A4, which restates the theorem-like environments the retained text needs and adds the article's own macro definitions that the retained text needs (\texttt{\textbackslash beas}, \texttt{\textbackslash eeas}, \texttt{\textbackslash epsilon}, \texttt{\textbackslash ol}), with extra wrapper packages (bm, bbm, dsfont, xspace, cancel, mathtools). The delivered numbering is the unit's own. Assets: no retained span contains a figure environment, a graphics-inclusion command, a PDF-page inclusion, a table or a TikZ picture; no image file is copied into this package, the package declares no asset dependency and no image file is redistributed with it. Licence: version arXiv:2511.05076v1 of this article is distributed under the Creative Commons Attribution 4.0 International licence, as recorded in the arXiv licence metadata for this exact version and shown on the abstract page https://arxiv.org/abs/2511.05076v1. A full rights notice is delivered at licenses/arxiv-2511.05076-CC-BY-4.0.txt. Characters: every retained excerpt is pure ASCII, so the delivered engine, XeLaTeX with its default Latin Modern text font, reports no missing character.
\begin{thm}\label{Th1}
Let $f=h\overline{g}$ be a non-vanishing sense-preserving and locally univalent log-harmonic mapping with the dilatation $\omega=g'h/h'g$ defined in $\mathbb{D}$. If $\log{g(z)}$ is an analytic Bloch function in $\mathbb{D}$, then either, $\Vert P_f \Vert =\Vert P_{hg^{\epsilon}}\Vert =\infty$ or, both $\Vert P_f \Vert $ and $\Vert P_{hg^{\epsilon}}\Vert $ are finite for each $\epsilon\in\ol{\mathbb{D}}$. If $\Vert P_f \Vert $ is finite, then
\beas
\left|\Vert P_f \Vert -\Vert P_{hg^{\epsilon}}\Vert \right|\leq 1+|1-\epsilon|\;\beta_{\log{g}}\leq 1+2\;\beta_{\log{g}},
\eeas
for each $\epsilon\in\ol{\mathbb{D}}$. 
In particular,
\beas
\left|\Vert P_f \Vert -\Vert P_{h}\Vert \right|\leq 1+\beta_{\log{g}}.\eeas
\end{thm}

\begin{thm}\label{Th11}
Let $f=h\overline{g}$ be a sense-preserving log-harmonic mapping such that $h$ is analytic and locally univalent, and $g$ is non-vanishing analytic in $\mathbb{D}$ with the dilatation $\omega=g'h/h'g$ defined in $\mathbb{D}$. If $\log{g(z)}$ is an analytic Bloch function in $\mathbb{D}$, then either, $\Vert P_f \Vert =\Vert P_{hg^{\epsilon}}\Vert =\infty$ or, both $\Vert P_f \Vert $ and $\Vert P_{hg^{\epsilon}}\Vert $ are finite for each $\epsilon\in\ol{\mathbb{D}}$. If $\Vert P_f \Vert $ is finite, then
\beas
\left|\Vert P_f \Vert -\Vert P_{hg^{\epsilon}}\Vert \right|\leq 1+2\;\beta_{\log{g}},
\eeas
for each $\epsilon\in\ol{\mathbb{D}}$. 
In particular,
\beas
\left|\Vert P_f \Vert -\Vert P_{h}\Vert \right|\leq 1+\beta_{\log{g}}.\eeas
Both the constants $ 1+2\;\beta_{\log{g}}$ and $1+\beta_{\log{g}}$ are sharp.
\end{thm}

\begin{proof} Using analogous reasoning to that used to prove \textrm{Theorem} \ref{Th1}, we arrive at the desired conclusion.\\[2mm]
\indent To show that the constant $1+2\beta_{\log g}$ is sharp, we consider the log-harmonic mapping $f(z)=h(z)\overline{g(z)}$ where $h(z)=z/(1-z)$, $g(z)=1/(1-z)$. It is 
evident that $h(z)$ is an analytic and locally univalent function while $g(z)$ is non-vanishing analytic in $\mathbb{D}$. As $g'(z)h(z)/(g(z)h'(z))=\omega(z)$, thus, we have the 
dilatation $\omega(z)=z$. Let $\epsilon=-1$. The pre-Schwarzian derivatives of $f$ and $hg^{-1}$ are, respectively, given by 
\beas
P_f(z)=\frac{3}{1-z}-\frac{\overline{z}}{1-|z|^2}\quad\text{and}\quad P_{hg^{-1}}(z)=0.\eeas
It is evident that
\beas
\beta_{\log{g}}=\sup_{z\in\mathbb{D}}(1-|z|^2)\left|\frac{g'(z)}{g(z)}\right|=\sup_{z\in\mathbb{D}}(1-|z|^2)\left|\frac{1}{1-z}\right|=2,\eeas
which shows that $\log{g}$ is analytic Bloch function. Thus, the pre-Schwarzian norms of $f$ and $hg^{-1}$ are, respectively, given by 
\beas
\Vert P_f\Vert=5\quad\text{and}\quad \Vert P_{hg^{-1}}\Vert=0.\eeas
Therefore, $\left|\Vert P_f\Vert-\Vert P_{hg^{-1}}\Vert\right|=5=1+2\cdot2$.\\[2mm]
\indent To show that the constant $1+\beta_{\log g}$ is sharp, we consider the log-harmonic mapping $f(z)=h(z)\overline{g(z)}$ with $h(z)=1/(1-z)$ and dilatation $\omega(z)=(a-z)/(1-az)$, where $a\in(0,1)$. As $\omega=g'h/h'g$, we have 
\beas \frac{g'(z)}{g(z)}=\frac{a-z}{(1-az)(1-z)}.\eeas
It is evident that
\beas
\beta_{\log{g}}=\sup_{z\in\mathbb{D}}(1-|z|^2)\left|\frac{g'(z)}{g(z)}\right|=\sup_{z\in\mathbb{D}}(1-|z|^2)\left|\frac{a-z}{(1-az)(1-z)}\right|.\eeas
On the positive real axis, we have 
\beas \sup_{0\leq r<1}\frac{(a-r)(1+r)}{(1-ar)}=2.\eeas
Therefore, $\log{g}$ is analytic Bloch function and $\beta_{\log{g}}=2$. It is evident that 
\beas 1-|\omega(z)|^2&=&\frac{(1-a^2)(1-|z|^2)}{(1-az)(1-a\ol{z})}.\eeas
Therefore, the pre-Schwarzian derivative of $f$ is 
\beas P_f(z)&=&\frac{h''(z)}{h'(z)}+\frac{g'(z)}{g(z)}-\frac{\overline{\omega(z)}\omega'(z)}{1-|\omega(z)|^2}\\[2mm]
%&=&\frac{2}{1-z}+\frac{a-z}{(1-az)(1-z)}-\frac{(a-\ol{z})}{(1-a\ol{z})}\cdot\frac{(a^2-1)}{(1-az)^2}\cdot\frac{(1-az)(1-a\ol{z})}{(1-a^2)(1-|z|^2)}\\[2mm]
&=&\frac{2}{1-z}+\frac{a-z}{(1-az)(1-z)}+\frac{(a-\ol{z})}{(1-az)(1-|z|^2)}\eeas
Similarly, a direct computation shows that $P_h(z)=h''(z)/h'(z)=2/(1-z)$ and the pre-Schwarzian norm of $h$ is $\Vert P_h\Vert=4$. 
The pre-Schwarzian norm of $f$ is 
\beas\Vert P_f\Vert=\sup_{z\in\mathbb{D}}\;(1-|z|^2)\left|\frac{2}{1-z}+\frac{a-z}{(1-az)(1-z)}+\frac{a-\ol{z}}{(1-az)(1-|z|^2)}\right|.\eeas
On the positive real axis, we have 
\beas \sup_{0\leq r<1}\left|2(1+r)+\frac{(a-r)(1+r)}{1-a r}+\frac{a-r}{1-a r}\right|=\sup_{0\leq r<1}G(a, r),\eeas
where $G(a,r)=2(1+r)+(a-r)(2+r)/(1-a r)$. Differentiate $G(a, r)$ partially with respect to $r$, we obtain
\beas \frac{\partial }{\partial r}G(a, r)=\frac{\left(2 a+1\right)\left(a r^2-2 r+a\right)}{(1-a r)^2}.\eeas
The roots of the equation $a r^2-2 r+a=0$ are 
\beas r_1(a)=\frac{1-\sqrt{1-a^2}}{a}\quad\text{and}\quad r_2(a)=\frac{1+\sqrt{1-a^2}}{a}.\eeas
It is evident that $r_2(a)>1$ and $r_1(a)\in(0, 1)$ for $a\in(0, 1)$. Hence, we have  
\beas\Vert P_f\Vert=G(a, r_1(a))=\frac{4 a^3+\left(\sqrt{1-a^2}+2\right) a^2+(4a+2) \left(\sqrt{1-a^2}-1\right)}{a^2 \sqrt{1-a^2}},\eeas
which tends to $7$ as $a\to1^-$. Therefore, we see that 
\beas \left|\Vert P_f \Vert -\Vert P_{h}\Vert \right|=3= 1+\beta_{\log{g}}.\eeas
This completes the proof.
\end{proof}

\end{document}
